\chapter{Free Groups and Covering Graphs}
\label{ch:viii14-free-groups-covering-graphs}

Graphs are the combinatorial laboratory of covering-space theory.  Their
fundamental groups are free, their universal covers are trees, and connected
covering graphs correspond to subgroups of free groups.  This chapter develops
the algebra of reduced words together with the geometry of graph covers.

\section{Free groups}

\begin{definition}[Free group]
\label{def:viii14-free-group}
Let \(S\) be a set.  The free group \(F(S)\) is the group of reduced finite
words in symbols
\[
s^{\pm1},\qquad s\in S,
\]
with multiplication given by concatenation followed by cancellation of adjacent
inverse pairs.
\end{definition}

If \(S=\{x_1,\dots,x_r\}\), write
\[
F_r=F(x_1,\dots,x_r).
\]

\section{Reduced words}

A word is reduced if it contains no adjacent subword
\[
xx^{-1}
\quad\text{or}\quad
x^{-1}x.
\]

\begin{proposition}[Unique reduced form]
\label{prop:viii14-reduced}
Every element of a free group has a unique reduced-word representative.
\end{proposition}

This normal form is the basis for concrete computations.

\section{Universal property}

\begin{theorem}[Universal property of a free group]
\label{thm:viii14-universal-free}
Given a set map
\[
f:S\to G
\]
into a group \(G\), there exists a unique homomorphism
\[
\overline f:F(S)\to G
\]
extending \(f\).
\end{theorem}

Thus a homomorphism out of a free group is completely determined by the images
of its free generators.

\section{Graphs as one-dimensional CW complexes}

A connected graph \(X\) is a one-dimensional CW complex.  Choose a base vertex
\(v_0\).  Edge paths determine loops and homotopies can be simplified
combinatorially by canceling immediate backtracking.

\section{Spanning trees}

\begin{definition}[Spanning tree]
\label{def:viii14-spanning-tree}
A spanning tree \(T\subset X\) of a connected graph is a connected subgraph
containing every vertex and having no cycles.
\end{definition}

A finite connected graph always has a spanning tree.

\section{Collapsing a spanning tree}

\begin{theorem}[Graph collapse]
\label{thm:viii14-collapse-tree}
If \(T\) is a spanning tree of a connected graph \(X\), then the quotient
\[
X/T
\]
is a wedge of circles, one for each edge of \(X\setminus T\), and
\[
X\simeq X/T.
\]
\end{theorem}

Hence graph fundamental groups are free.

\section{Rank formula}

For a finite connected graph with \(V\) vertices and \(E\) edges, a spanning
tree has \(V-1\) edges.  Therefore the number of edges outside the tree is
\[
E-(V-1)=E-V+1.
\]

\begin{theorem}[Fundamental group of a finite connected graph]
\label{thm:viii14-rank}
\[
\pi_1(X)\cong F_{E-V+1}.
\]
\end{theorem}

\section{Wedges of circles}

For the bouquet
\[
\bigvee_{j=1}^rS^1,
\]
one obtains
\[
\pi_1\left(\bigvee_{j=1}^rS^1\right)\cong F_r.
\]

Each circle contributes one free generator.

\section{Trees}

\begin{definition}[Tree]
\label{def:viii14-tree}
A tree is a connected graph containing no nontrivial cycles.
\end{definition}

\begin{theorem}[Trees are simply connected]
\label{thm:viii14-tree-simply}
Every tree is contractible and therefore simply connected.
\end{theorem}

A finite tree collapses edge by edge to a vertex.  Infinite trees admit the
same path-uniqueness intuition and are simply connected as graphs.

\section{Universal covers of graphs}

\begin{theorem}[Universal covering graph]
\label{thm:viii14-universal-tree}
The universal cover of a connected graph is a tree.
\end{theorem}

Conversely, a tree with a free properly discontinuous graph automorphism action
gives a graph quotient whose universal cover is that tree.

\section{Cayley graphs}

For a group \(G\) generated by a set \(S\), the Cayley graph has vertex set
\(G\) and directed edges
\[
g\longrightarrow gs
\]
for \(s\in S\).

\begin{example}[Cayley tree of a free group]
\label{ex:viii14-cayley-tree}
The Cayley graph of \(F_r\) with respect to its free generators is a regular
tree of valence \(2r\).
\end{example}

\section{Deck action of a free group}

The universal cover of the wedge of \(r\) circles is the Cayley tree of
\(F_r\).  The group acts by left multiplication on vertices and edges.

\begin{theorem}[Free-group deck action]
\label{thm:viii14-deck-free}
\[
F_r
\cong
\operatorname{Deck}(\widetilde X\to X)
\]
for the universal cover of the \(r\)-circle wedge.
\end{theorem}

\section{Covering graphs}

A graph map
\[
p:\widetilde X\to X
\]
is a covering graph when the incident edge pattern around every vertex maps
bijectively to the incident edge pattern around its image.  This is the
combinatorial form of an evenly covered neighborhood.

\section{Subgroups and connected graph covers}

For a connected based graph \(X\), pointed connected covers correspond to
subgroups
\[
H\le\pi_1(X).
\]

Since \(\pi_1(X)\) is free, graph coverings give a geometric realization of
subgroups of free groups.

\section{Schreier graphs}

Let \(F_r\) act on the right cosets of a subgroup \(H\).  The associated
Schreier graph has vertices
\[
Hg,\qquad g\in F_r,
\]
and edges labeled by free generators.

This graph realizes the covering corresponding to \(H\).

\section{Nielsen--Schreier theorem}

\begin{theorem}[Nielsen--Schreier]
\label{thm:viii14-nielsen-schreier}
Every subgroup of a free group is free.
\end{theorem}

A covering-space proof is immediate in spirit: a subgroup determines a
connected covering graph, and the fundamental group of that graph is free.

\section{Finite-index rank formula}

\begin{theorem}[Schreier index formula]
\label{thm:viii14-schreier-formula}
If
\[
H\le F_r
\]
has finite index \(d\), then
\[
\operatorname{rank}(H)=1+d(r-1).
\]
\end{theorem}

This can be seen by counting vertices and edges in the finite covering graph.

\section{Regular graph covers}

A covering graph corresponding to
\[
H\le F_r
\]
is regular exactly when \(H\) is normal.  Then the deck group is
\[
F_r/H.
\]

\section{Core graphs}

For finitely generated subgroups of free groups, one can prune hanging trees
from a covering graph and retain a finite core containing every reduced loop.
This gives a compact combinatorial model of the subgroup.

\section{Edge paths and reduced words}

In a rose with labeled oriented edges
\[
x_1,\dots,x_r,
\]
a based edge loop reads a word in the symbols \(x_i^{\pm1}\).  Removing
immediate backtracking corresponds exactly to free reduction.

Thus graph homotopy and free-group reduction are two versions of the same
operation.

\section{Bridge to simplicial complexes}

Graphs are one-dimensional simplicial complexes after subdividing loops and
multiple edges if necessary.  VIII/15 generalizes the vertex-edge language to
higher-dimensional simplices and prepares the chain complexes used in homology.


\section{Exercises}

\begin{exercise}\label{exr:viii14-01}
Define \(F(S)\).
\end{exercise}
\begin{hint}
Use reduced words.
\end{hint}
\begin{solution}
The group of reduced finite words in \(s^{\pm1}\) for \(s\in S\).
\end{solution}

\begin{exercise}\label{exr:viii14-02}
What is a reduced word?
\end{exercise}
\begin{hint}
No adjacent inverse pair.
\end{hint}
\begin{solution}
A word containing neither \(xx^{-1}\) nor \(x^{-1}x\) adjacently.
\end{solution}

\begin{exercise}\label{exr:viii14-03}
What is \(F_r\)?
\end{exercise}
\begin{hint}
Choose \(r\) free generators.
\end{hint}
\begin{solution}
The free group on \(r\) generators.
\end{solution}

\begin{exercise}\label{exr:viii14-04}
State the universal property of \(F(S)\).
\end{exercise}
\begin{hint}
Extend any set map to a homomorphism.
\end{hint}
\begin{solution}
Every map \(S\to G\) extends uniquely to \(F(S)\to G\).
\end{solution}

\begin{exercise}\label{exr:viii14-05}
Define a spanning tree.
\end{exercise}
\begin{hint}
All vertices, connected, no cycles.
\end{hint}
\begin{solution}
A connected cycle-free subgraph containing every vertex.
\end{solution}

\begin{exercise}\label{exr:viii14-06}
How many edges does a finite spanning tree on \(V\) vertices have?
\end{exercise}
\begin{hint}
Use the tree formula.
\end{hint}
\begin{solution}
\(V-1\).
\end{solution}

\begin{exercise}\label{exr:viii14-07}
What is \(X/T\) for a spanning tree \(T\)?
\end{exercise}
\begin{hint}
Collapse the tree.
\end{hint}
\begin{solution}
A wedge of circles indexed by edges outside \(T\).
\end{solution}

\begin{exercise}\label{exr:viii14-08}
Give the rank of \(\pi_1(X)\) for a finite connected graph.
\end{exercise}
\begin{hint}
Count non-tree edges.
\end{hint}
\begin{solution}
\(E-V+1\).
\end{solution}

\begin{exercise}\label{exr:viii14-09}
What is \(\pi_1(\bigvee^rS^1)\)?
\end{exercise}
\begin{hint}
Each circle gives one generator.
\end{hint}
\begin{solution}
\(F_r\).
\end{solution}

\begin{exercise}\label{exr:viii14-10}
Define a tree.
\end{exercise}
\begin{hint}
Connected graph with no cycles.
\end{hint}
\begin{solution}
A connected cycle-free graph.
\end{solution}

\begin{exercise}\label{exr:viii14-11}
What is \(\pi_1\) of a tree?
\end{exercise}
\begin{hint}
Trees are contractible.
\end{hint}
\begin{solution}
Trivial.
\end{solution}

\begin{exercise}\label{exr:viii14-12}
What is the universal cover of a connected graph?
\end{exercise}
\begin{hint}
Use simple connectivity.
\end{hint}
\begin{solution}
A tree.
\end{solution}

\begin{exercise}\label{exr:viii14-13}
What is the Cayley graph vertex set of \(G\)?
\end{exercise}
\begin{hint}
One vertex per group element.
\end{hint}
\begin{solution}
\(G\).
\end{solution}

\begin{exercise}\label{exr:viii14-14}
What valence has the standard Cayley graph of \(F_r\)?
\end{exercise}
\begin{hint}
Each generator and inverse gives an edge direction.
\end{hint}
\begin{solution}
\(2r\).
\end{solution}

\begin{exercise}\label{exr:viii14-15}
What group acts by deck transformations on the universal cover of the \(r\)-rose?
\end{exercise}
\begin{hint}
Use the fundamental group.
\end{hint}
\begin{solution}
\(F_r\).
\end{solution}

\begin{exercise}\label{exr:viii14-16}
What local condition defines a covering graph?
\end{exercise}
\begin{hint}
Incident edges must match bijectively.
\end{hint}
\begin{solution}
At every vertex, the star maps bijectively to the star of its image.
\end{solution}

\begin{exercise}\label{exr:viii14-17}
What algebraic objects classify pointed connected graph covers?
\end{exercise}
\begin{hint}
Use covering classification.
\end{hint}
\begin{solution}
Subgroups of the graph fundamental group.
\end{solution}

\begin{exercise}\label{exr:viii14-18}
What is a Schreier graph?
\end{exercise}
\begin{hint}
Use cosets of a subgroup.
\end{hint}
\begin{solution}
A labeled graph whose vertices are cosets \(Hg\) and whose edges encode multiplication by free generators.
\end{solution}

\begin{exercise}\label{exr:viii14-19}
State Nielsen--Schreier.
\end{exercise}
\begin{hint}
Subgroups preserve freeness.
\end{hint}
\begin{solution}
Every subgroup of a free group is free.
\end{solution}

\begin{exercise}\label{exr:viii14-20}
State Schreier's finite-index rank formula.
\end{exercise}
\begin{hint}
Use index \(d\).
\end{hint}
\begin{solution}
\(\operatorname{rank}(H)=1+d(r-1)\).
\end{solution}

\begin{exercise}\label{exr:viii14-21}
When is a connected graph cover regular?
\end{exercise}
\begin{hint}
Use subgroup normality.
\end{hint}
\begin{solution}
When the corresponding subgroup is normal.
\end{solution}

\begin{exercise}\label{exr:viii14-22}
What is its deck group in the normal case?
\end{exercise}
\begin{hint}
Take quotient.
\end{hint}
\begin{solution}
\(F_r/H\).
\end{solution}

\begin{exercise}\label{exr:viii14-23}
What does free reduction correspond to geometrically?
\end{exercise}
\begin{hint}
Cancel immediate edge backtracking.
\end{hint}
\begin{solution}
Removing an edge immediately followed by the same edge in reverse.
\end{solution}

\begin{exercise}\label{exr:viii14-24}
What is the bridge to VIII/15?
\end{exercise}
\begin{hint}
Graphs are 1-dimensional simplicial objects.
\end{hint}
\begin{solution}
VIII/15 generalizes vertices and edges to higher-dimensional simplices.
\end{solution}


\section{Solved problem dossiers}

\begin{problem}[Reduce a word]\label{prob:viii14-01}
Reduce \(aba^{-1}ab^{-1}b\) in the free group \(F(a,b)\).
\end{problem}
\begin{solution}
The terminal \(b^{-1}b\) cancels, leaving \(aba^{-1}a\). No further adjacent inverse pair occurs.
\end{solution}

\begin{problem}[Free universal property]\label{prob:viii14-02}
Define a homomorphism \(F(a,b)\to\mathbb Z\) by \(a\mapsto2\), \(b\mapsto-1\). What is the image of \(ab^{-1}a\)?
\end{problem}
\begin{solution}
By the universal property, extend additively: \(2-(-1)+2=5\).
\end{solution}

\begin{problem}[Triangle graph rank]\label{prob:viii14-03}
Compute the fundamental-group rank of a triangle graph.
\end{problem}
\begin{solution}
It has \(V=3\), \(E=3\), so the rank is \(3-3+1=1\), hence \(\pi_1\cong\mathbb Z\).
\end{solution}

\begin{problem}[Square with diagonal]\label{prob:viii14-04}
Compute the rank for a square with one diagonal.
\end{problem}
\begin{solution}
Here \(V=4\), \(E=5\), so rank \(=5-4+1=2\).
\end{solution}

\begin{problem}[Complete graph \(K_4\)]\label{prob:viii14-05}
Compute the rank of \(\pi_1(K_4)\).
\end{problem}
\begin{solution}
\(K_4\) has \(V=4\), \(E=6\), so the rank is \(3\), hence \(\pi_1(K_4)\cong F_3\).
\end{solution}

\begin{problem}[Tree criterion]\label{prob:viii14-06}
Show a finite connected graph is a tree iff \(E=V-1\).
\end{problem}
\begin{solution}
A spanning tree already has \(V-1\) edges. If the connected graph has no extra edges it is the tree; any extra edge creates a cycle. Conversely every tree has \(V-1\) edges.
\end{solution}

\begin{problem}[Rose cover]\label{prob:viii14-07}
Describe the universal cover of \(S^1\vee S^1\).
\end{problem}
\begin{solution}
It is the 4-regular Cayley tree of \(F_2\), with edges labeled by the two generators and their inverse orientations.
\end{solution}

\begin{problem}[Deck action]\label{prob:viii14-08}
How does \(F_2\) act on that tree?
\end{problem}
\begin{solution}
By left multiplication on vertex labels \(g\mapsto hg\), carrying each labeled edge to the corresponding labeled edge.
\end{solution}

\begin{problem}[Index-two subgroup rank]\label{prob:viii14-09}
If \(H\le F_3\) has index \(2\), find its rank.
\end{problem}
\begin{solution}
Schreier gives \(1+2(3-1)=5\).
\end{solution}

\begin{problem}[Index-five subgroup of \(F_2\)]\label{prob:viii14-10}
Find the rank.
\end{problem}
\begin{solution}
\(1+5(2-1)=6\).
\end{solution}

\begin{problem}[Subgroup of \(\mathbb Z\)]\label{prob:viii14-11}
Interpret \(n\mathbb Z\le F_1\cong\mathbb Z\) as a graph cover.
\end{problem}
\begin{solution}
It corresponds to the \(n\)-sheeted circle cover \(S^1\to S^1\), with covering graph an \(n\)-cycle.
\end{solution}

\begin{problem}[Nielsen--Schreier proof idea]\label{prob:viii14-12}
Explain the covering-space proof that every subgroup \(H\le F_r\) is free.
\end{problem}
\begin{solution}
Realize \(F_r\) as \(\pi_1\) of the \(r\)-rose. The subgroup gives a connected covering graph. Fundamental groups of graphs are free, and that group identifies with \(H\).
\end{solution}

\begin{problem}[Normal subgroup cover]\label{prob:viii14-13}
What extra symmetry appears when \(H\trianglelefteq F_r\)?
\end{problem}
\begin{solution}
The associated cover is regular, and deck transformations act transitively on fibers with group \(F_r/H\).
\end{solution}

\begin{problem}[Non-normal subgroup]\label{prob:viii14-14}
What does the deck group become for a general subgroup?
\end{problem}
\begin{solution}
Under the covering classification, it is \(N_{F_r}(H)/H\), which can be smaller than the full quotient of \(F_r\).
\end{solution}

\begin{problem}[Schreier graph vertices]\label{prob:viii14-15}
If \(H\le F_r\) has index \(d\), how many vertices does its Schreier graph have?
\end{problem}
\begin{solution}
Exactly \(d\), one for each right coset.
\end{solution}

\begin{problem}[Schreier edge count]\label{prob:viii14-16}
For a degree-\(d\) cover of the \(r\)-rose, how many unoriented edges are there?
\end{problem}
\begin{solution}
Each of the \(r\) base loops lifts with one outgoing edge from each of \(d\) vertices, giving \(rd\) unoriented labeled edges.
\end{solution}

\begin{problem}[Derive Schreier rank]\label{prob:viii14-17}
Use the previous count to recover the index formula.
\end{problem}
\begin{solution}
The connected covering graph has \(V=d\), \(E=rd\), hence rank \(E-V+1=rd-d+1=1+d(r-1)\).
\end{solution}

\begin{problem}[Core graph]\label{prob:viii14-18}
Why can hanging trees be removed without changing the subgroup represented by a covering graph?
\end{problem}
\begin{solution}
Hanging trees contain no reduced closed loops. Collapsing or pruning them does not alter the fundamental group, leaving the loop-carrying core.
\end{solution}

\begin{problem}[Backtracking cancellation]\label{prob:viii14-19}
Why does an edge followed immediately by its reverse represent a null segment in a graph path?
\end{problem}
\begin{solution}
That out-and-back segment contracts through the edge, so deleting it is an endpoint-fixing path homotopy.
\end{solution}

\begin{problem}[Bridge to simplices]\label{prob:viii14-20}
How does a graph become a simplicial complex?
\end{problem}
\begin{solution}
Use vertices as 0-simplices and non-loop edges as 1-simplices; subdivide loops or multiple-edge configurations as needed to fit the abstract simplicial-complex rules.
\end{solution}
